已有工作负载研究说明，Poisson 近似应由具体时间尺度上的离散度与相关性诊断决定\cite{zou_poissonity_2022}；分层需求预测则强调在细粒度结构与参数稳定性之间进行数据驱动聚合\cite{cohen_dac_2022}。本文据此保留训练段识别出的联合业务标记和类型条件资源标记。

\begin{algorithm}[H]
\SetAlgoLined
\setstretch{0.82}
\SetAlFnt{\footnotesize}
\SetArgSty{textnormal}
\SetKwInput{KwIn}{输入}
\SetKwInput{KwOut}{输出}
\SetKwFor{ForEach}{对每个}{执行}{结束}
\SetKwIF{If}{ElseIf}{Else}{若}{则}{否则若}{否则}{结束}
\SetKwFunction{Diag}{结构诊断}
\SetKwFunction{Fit}{拟合MCP}
\SetKwFunction{MC}{蒙特卡洛预测}
\SetKwFunction{Place}{本地优先放置}
\caption{两阶段短期概率预测与基础调度}
\label{alg:q1-unified}
\KwIn{$\mathcal D_{0:2399}$，$\mathcal J_{2376:2399}$，容量与 SLA，$M=5000$}
\KwOut{24 h 概率预测、实际任务排程与约束审计}

\tcp*[l]{阶段 I：结构识别与独立验证}
$(\mathcal D_{tr},\mathcal D_{val},\mathcal D_{te})\leftarrow(0{:}2351,2352{:}2375,2376{:}2399)$\;
$\mathcal M\leftarrow\Diag(\mathcal D_{tr})$ \tcp*[r]{Fano/ACF、区域--类型关联、$(G,D)\mid K$}
$\theta_1=(\hat\lambda,\hat\pi_{rk},\widehat F_{GD}^{(k)})\leftarrow\Fit(\mathcal M,\mathcal D_{tr})$\;
$\mathcal I_{val}\leftarrow\MC(\theta_1,M)$，计算 MAE、WAPE、RMSE 与 PICP\;
验证通过后冻结 $\mathcal M$，测试段不参与模型选择\;

\tcp*[l]{阶段 II：参数重估、最终预测与基础调度}
$\theta_2\leftarrow\Fit(\mathcal M,\mathcal D_{0:2375})$，$\mathcal I_{te}\leftarrow\MC(\theta_2,M)$\;
初始化 2376 h 前未完成任务形成的 GPU/IT/Facility 占用\;
\ForEach{按到达顺序的 $j\in\mathcal J_{2376:2399}$}{
  构造合法域 $C_j=\{(r,s):\mathrm{SLA},\mathrm{LatestFinish},s+p_j\le2406\}$；RT 固定 $s=a_j$\;
  $(r_j^*,s_j^*)\leftarrow\Place(C_j)$，同时满足 GPU/IT/Facility 容量\;
  \If{$(r_j^*,s_j^*)=\varnothing$}{返回不可行排程\;}
  更新区间 $[s_j^*,s_j^*+p_j)$ 的资源占用\;
}
审计唯一执行、RT 即时开工、容量、SLA、LatestFinish 与 2406 边界\;
\end{algorithm}
